\chapter{Tonal Processes}
\label{chap:process}

\section{Overview}
\label{sec:process-overview}

This chapter turns to tonal processes, that is, the mapping
between the underlying forms (UR) and the surface forms (SRs). 
In particular, I focus mainly on those processes that feature 
associations lines changes, for it is these processes that 
typically present the most challenges for learning. The goal of 
this chapter is to demonstrate that these processes can be learned 
using logical transductions.

As Chapter~\ref{chap:tonotactics}\todo{fix the structure of prev chap}, 
I will first present relevant linguistic data for three types of 
associational changes: processes triggered by segmental changes, 
processes triggered by tone quality changes, and tone-segment 
interactions (\ref{sec:process-type}). I then formalize the learning 
problem and present the logical transduction framework (\ref{sec:logi_trans}, some background on logical transductions is in 
Chapter~\ref{chap:formal}). I then propose a learning algorithm and 
demonstrate the learning process using three case studies.

\section{Tonal Association Changes}
\label{sec:process-type}
In this section, I present three types of tonal processes that
primarily involve changes in association lines. As shown in
Table~\ref{tab:process-types}, these processes fall into three groups
according to the source of the trigger: (i)~segmentally motivated
changes, (ii)~changes motivated by tone quality, and
(iii)~changes arising from tone--segment interaction. All three types
reflect the autonomy of tones from their host TBUs: when a change 
occurs on one tier, the associations between tiers are rearranged to
accommodate the change. Type~(i) is triggered by a change on the TBU tier,
type~(ii) by a change on the tone tier, and type~(iii) by the
interaction of both tiers.

\begin{table}[htbp]
  \centering
  \caption{Tone reassociation processes classified by trigger.}
  \label{tab:process-types}
  \renewcommand{\arraystretch}{1.2}
  \begin{tabularx}{\textwidth}{@{} >{\RaggedRight}p{0.2\textwidth}
                                  >{\RaggedRight}X
                                  >{\RaggedRight}p{0.24\textwidth} @{}}
    \toprule
    Trigger & Description & Example \\
    \midrule
    Segment-driven
      & reassociate to accommodate TBU deletion/insertion, or syllabic shape
      & Etako, Engenni, Ogoja Yala \\ 
    Tone-driven
      & reassociate to accommodate a particular melody
      & Spreading, OCP \\
    Tone--TBU
      & reassociate to accommodate a tone-TBU context
      & Consonant depressors \\
    \bottomrule
  \end{tabularx}
\end{table}

As we progress, one can see that all three types of processes can be captured through a single association relation
\(\alpha\) together with two sub-operations, \textsf{deassoc} and
\textsf{reassoc}, modeled as logical transductions, and be learned by
a BUFIA-styled algorithm as shown in Chapter\ref{chap:learning}. 
There are indeed processes that the current learning fails to 
capture yet, and I will discuss them in the section \todo{add 
label}. 




\subsection{Segmentally Motivated Changes}

One significant consequence of ARs is it is able to capture
the tonal stability. One common case is triggered by vowel 
deletion, which is often accompanied by tonal changes. Here,
I use vowel deletion as a remedy for vowel hiatus from three
languages: Engenni, Estako, and Ogoja Yala. When in vowel
hiatus, all three languages will delete the first vowel and
keep the following vowel. However, the tonal behaviors are
differ.

\ex \textbf{Engenni} \label{ex:engenni}
\begin{multicols}{2}
\quad\ L\quad H\qquad\qquad H\\
 /\` ɔkò \# édèì/ $\rightarrow$ [\` ɔkédèì] \\
 \textit{man \ \ canoe} $\rightarrow$ \textit{``a man’s canoe"} \\
 
 \quad\ \ H\quad L\qquad\qquad\ \ \ L\\
/fùnú \# èdhí/ $\rightarrow$ [fùnèdhí] \\
\textit{climb \ palm tree} $\rightarrow$ \textit{``to climb a palm tree"}
\end{multicols}
\xe

In (\ref{ex:engenni}), Engenni deletes the first vowel in vowel hiatus. When the vowel is deleted, the tone of the deleted vowel is also deleted. The surviving vowel keeps its original tone. To represent this process autosegmentally, the only change regarding the association line change is a delinking between T$_1$ and V$_1$. The other association lines stay the same.
The AR is shown in (\ref{ex:engenni-etsako}a).

Another example is from Etsako (Example \ref{ex:etsako}), which also deletes the
first vowel in the hiatus, but the tone of the deleted
vowel will survive the deletion, and forms a new link 
to the following vowel that survives. Since the surviving vowel also bears its original tone, the two
tone will form a contour (following the original order) on the surviving vowel. 


\ex \textbf{Etsako (contour formation):} 
\label{ex:etsako}
\vspace{-1em}
\begin{multicols}{2}
\quad\ \ L \ \ \ H\qquad\qquad LH\\
/ówà \# ówà/ \ $\rightarrow$ [ówǒwà] \\
\textit{house \ \ house} $\rightarrow$ \textit{``every house"}

\qquad\ H \ \ \ L\qquad\qquad HL\\
/\`ɔ \# dé \ \# àkpà/ $\rightarrow$ [\`ɔdâkpà] \\
\textit{he \ \ buy \ \ cup}\quad$\rightarrow$ \textit{``he buys a cup"}
\end{multicols}
\xe


Similar to Engenni, Estako also has a delink takes place between T$_1$ and V$_1$. However, the tone of the deleted vowel (T$_1$) will be reassociated to the surviving vowel (V$_2$). The AR is shown in (\ref{ex:engenni-etsako}b).


\ex \label{ex:engenni-etsako}
  a.~\begin{tikzpicture}[scale=1.2, baseline=(t1.base)]
      \node[tone] (t1) at (0,1) {T$_1$};
      \node[tone] (t2) at (1,1) {T$_2$};
      \node[seg]  (v1) at (0,0) {V$_1$};
      \node[seg]  (v2) at (1,0) {V$_2$};
      %%%% drawing
      \draw[assoc] (t1)--(v1) node[midway] {$=$};
      \draw[assoc] (t2)--(v2);
      \node (arrow) at (2,0.5) {$\rightarrow$};
      \node[tone] (t2o) at (3,1) {T$_2$};
      \node[seg]  (v2o) at (3,0) {V$_2$};
      \draw[assoc] (t2o)--(v2o);
    \end{tikzpicture}
  \hfill
  b.~\begin{tikzpicture}[scale=1.2, baseline=(t1.base)]
      \node[tone] (t1) at (0,1) {T$_1$};
      \node[tone] (t2) at (1,1) {T$_2$};
      \node[seg]  (v1) at (0,0) {V$_1$};
      \node[seg]  (v2) at (1,0) {V$_2$};
      %%%% drawing
      \draw[assoc] (t1)--(v1) node[midway] {$=$};
      \draw[assoc] (t2)--(v2);
      \draw[newline] (t1)--(v2);
      \node (arrow) at (2,0.5) {$\rightarrow$};
      \node[tone] (t1o) at (3,1) {T$_1$};
      \node[tone] (t2o) at (4,1) {T$_2$};
      \node[seg]  (v2o) at (3.5,0) {V$_2$};
      \draw[assoc] (t1o)--(v2o);
      \draw[assoc] (t2o)--(v2o);
    \end{tikzpicture}
  \hspace*{\fill}
\xe

The third case is from Ogoja Yala (Example \ref{ex:ogoja}), which 
always deletes the first vowel in the hiatus. However, whether the 
tone will be deleted or not depends on its pitch relative to its 
neighbor tone: if it is higher than following tone, T$_1$ will replace the T$_2$; but if it is lower, it will be deleted along with the deleting vowel. In addition, Ogoja Yala also has a mid tone (M). If no H is present (i.e., only L and M), then M is preferred. 

\ex Ogoja Yala
\label{ex:ogoja}
\begin{multicols}{3}
\quad L\quad H\qquad\ \ \ \quad H\\
/dè \# áchí/ $\rightarrow$ [dáchí]\\
\textit{give \ grass} $\rightarrow$ \textit{``to give grass"}

\quad H\quad L\qquad\ \ \ \quad H\\
/má \# \` ɔchí/ $\rightarrow$ [m\' ɔchí]\\
\vspace{1em}\textit{see \quad tree} $\rightarrow$ \textit{``to see tree"}

\hspace{0.2cm}M \ \ L\qquad\ \ \quad M\\
/b\=i \# \` ɔchí/ $\rightarrow$ [b\=ɔchí]\\
\textit{carry  stick} $\rightarrow$ \textit{``to carry a stick"}
\end{multicols}
\xe

The AR generalizations are shown in \ref{ex:ogoja-ar}. When T1 is L, it patterns with the Engenni case; when T1 is H, it patterns with the Etsako case. When T1 is M, the following tone determines the pattern.

\ex\label{ex:ogoja-ar}
  a.~\begin{tikzpicture}[scale=1.2, baseline=(t1.base)]
      \node[tone] (t1) at (0,1) {L};
      \node[tone] (t2) at (1,1) {\{H, M\}};
      \node[seg]  (v1) at (0,0) {V$_1$};
      \node[seg]  (v2) at (1,0) {V$_2$};
      %%%% drawing
      \draw[assoc] (t1)--(v1) node[midway] {$=$};
      \draw[assoc] (t2)--(v2);
      \node (arrow) at (2,0.5) {$\rightarrow$};
      \node[tone] (t2o) at (3,1) {\{H, M\}};
      \node[seg]  (v2o) at (3,0) {V$_2$};
      \draw[assoc] (t2o)--(v2o);
    \end{tikzpicture}
  \hfil
  b.~\begin{tikzpicture}[scale=1.2, baseline=(t1.base)]
      \node[tone] (t1) at (0,1) {H};
      \node[tone] (t2) at (1,1) {\{M,L\}};
      \node[seg]  (v1) at (0,0) {V$_1$};
      \node[seg]  (v2) at (1,0) {V$_2$};
      %%%% drawing
      \draw[assoc] (t1)--(v1) node[midway] {$=$};
      \draw[assoc] (t2)--(v2) node[midway] {$=$};
      \draw[newline] (t1)--(v2);
      \node (arrow) at (2,0.5) {$\rightarrow$};
      \node[tone] (t2o) at (3,1) {\{H, M\}};
      \node[seg]  (v2o) at (3,0) {V$_2$};
      \draw[assoc] (t2o)--(v2o);
      \draw[assoc] (t1o)--(v2o);
      \draw[assoc] (t2o)--(v2o);
    \end{tikzpicture}


    c.~\begin{tikzpicture}[scale=1.2, baseline=(t1.base)]
      \node[tone] (t1) at (0,1) {M};
      \node[tone] (t2) at (1,1) {H};
      \node[seg]  (v1) at (0,0) {V$_1$};
      \node[seg]  (v2) at (1,0) {V$_2$};
      %%%% drawing
      \draw[assoc] (t1)--(v1) node[midway] {$=$};
      \draw[assoc] (t2)--(v2);
      \node (arrow) at (2,0.5) {$\rightarrow$};
      \node[tone] (t2o) at (3,1) {H};
      \node[seg]  (v2o) at (3,0) {V$_2$};
      \draw[assoc] (t2o)--(v2o);
    \end{tikzpicture}
    \hfil
    d.~\begin{tikzpicture}[scale=1.2, baseline=(t1.base)]
      \node[tone] (t1) at (0,1) {M};
      \node[tone] (t2) at (1,1) {L};
      \node[seg]  (v1) at (0,0) {V$_1$};
      \node[seg]  (v2) at (1,0) {V$_2$};
      %%%% drawing
      \draw[assoc] (t1)--(v1) node[midway] {$=$};
      \draw[assoc] (t2)--(v2) node[midway] {$=$};
      \draw[newline] (t1)--(v2);
      \node (arrow) at (2,0.5) {$\rightarrow$};
      \node[tone] (t2o) at (3,1) {M};
      \node[seg]  (v2o) at (3,0) {V$_2$};
      \draw[assoc] (t2o)--(v2o);
      \draw[assoc] (t1o)--(v2o);
      \draw[assoc] (t2o)--(v2o);
    \end{tikzpicture}
    \hfil
  
\xe



\subsection{Tone-driven changes}
\subsection{Tone--segment interaction}


\section{Learning association changes using logical transductions}
As mentioned above, the three types under discussion can be captured through a single association relation
\(\alpha\) together with two sub-operations, \textsf{deassoc} and
\textsf{reassoc}, modeled as logical transductions. These
transductions are grounded in model-theoretic representations of
phonological structures, where logical formulas define the mapping
from input to output representations.
\section{The Learning Problem}

\section{The Learning Algorithm}

\section{Case Studies}

\section{Discussion}

\section{Conclusion}